\documentclass[aspectratio=169, 11pt]{beamer}

\usetheme{Madrid}
\usecolortheme{default}
\setbeamertemplate{navigation symbols}{}
\setbeamertemplate{footline}{
    \hfill
    \usebeamerfont{page number in head/foot}
    \insertframenumber{} / 14
    \hspace{0.4cm}
    \vspace{0.2cm}
}
\usepackage[T2A]{fontenc}
\usepackage[utf8]{inputenc}
\usepackage[russian]{babel}
\usepackage{booktabs}
\usepackage{tikz}
\usetikzlibrary{arrows.meta, positioning}

\title[Кольцевой буфер для диагностики ВМ]{
    Применение многопоточного кольцевого буфера для организации
    высокоскоростного обмена диагностической информацией в среде виртуализации
}
\author{Абубакиров Владимир Андреевич}
\institute{
    МФТИ, ФПМИ \\
    Кафедра системного программирования ИСП РАН \\
    \vspace{0.2cm}
    Научный руководитель: к.ф.-м.н. Камкин А. С. \\
    Научный консультант: Скородумов Д. П.
}
\date{Москва, 2026}

\begin{document}

\begin{frame}
    \titlepage
\end{frame}

\begin{frame}{Проблема}
    \begin{itemize}
        \item Современные виртуальные машины включают вспомогательные модули для более бесшовной работы.
        \item Такие модули отвечают за общий буфер обмена, USB-перенаправление, синхронизацию времени и разрешения экрана.
        \item Для отладки нужны диагностические сообщения из гостевой системы на хост.
        \item Важны низкая задержка, высокая пропускная способность и корректная работа нескольких источников.
        \item Отдельная сложность --- ранняя загрузка UEFI, когда общая память ещё недоступна.
    \end{itemize}
\end{frame}

\begin{frame}{Цель и задачи}
    \textbf{Цель:} разработать и оценить комплекс компонентов для высокоскоростной передачи диагностической информации в среде виртуализации.

    \vspace{0.4cm}
    \begin{itemize}
        \item Проанализировать существующие каналы передачи диагностических сообщений.
        \item Обосновать выбор MPSC-кольцевого буфера.
        \item Реализовать виртуальное устройство LogDev для Stratovirt.
        \item Реализовать поддержку UEFI и Windows.
        \item Проверить корректность прототипа и измерить накладные расходы.
    \end{itemize}
\end{frame}

\begin{frame}{Каналы передачи}
    \begin{table}
        \small
        \begin{tabular}{p{0.26\textwidth} p{0.28\textwidth} p{0.34\textwidth}}
            \toprule
            Подход & Преимущества & Ограничения \\
            \midrule
            Serial/debugcon &
            Простота, ранняя загрузка &
            Частые выходы в эмулятор при побайтовой записи \\
            \addlinespace
            virtio-serial &
            Стандартный виртуализированный канал &
            Требует инициализации драйвера и неудобен для ранних стадий UEFI \\
            \bottomrule
        \end{tabular}
    \end{table}
\end{frame}

\begin{frame}{Организация кольцевого буфера}
    \begin{table}
        \scriptsize
        \begin{tabular}{p{0.25\textwidth} p{0.31\textwidth} p{0.34\textwidth}}
            \toprule
            Подход & Преимущества & Ограничения \\
            \midrule
            SPSC &
            Простая и быстрая схема &
            Несколько писателей требуют внешнего мьютекса \\
            \addlinespace
            eBPF-style ring buffer &
            Общий буфер событий с reserve/copy/commit &
            Резервирование места защищено блокировкой \\
            \addlinespace
            MPSC с атомарным резервированием &
            Несколько писателей, резервирование через атомарные операции &
            Нужно не открывать читателю незавершённые записи \\
            \bottomrule
        \end{tabular}
    \end{table}
\end{frame}

\begin{frame}{Идея решения}
    \begin{columns}[T]
        \begin{column}{0.52\textwidth}
            \begin{itemize}
                \item Основной канал --- кольцевой буфер в RAM гостевой системы.
                \item Несколько писателей: UEFI-модули, приложения и драйверы Windows.
                \item Один читатель: виртуальное устройство LogDev в Stratovirt.
                \item MMIO используется для настройки и уведомлений, а не для передачи каждого байта.
            \end{itemize}
        \end{column}
        \begin{column}{0.45\textwidth}
            \centering
            \begin{tikzpicture}[
                node distance=0.55cm,
                box/.style={draw, rounded corners, align=center, minimum width=3.5cm, minimum height=0.8cm},
                arr/.style={-{Latex[length=2mm]}, thick}
            ]
                \node[box] (guest) {Гостевая система\\UEFI / Windows};
                \node[box, below=of guest] (buf) {Общая память\\MPSC ring buffer};
                \node[box, below=of buf] (dev) {LogDev\\Stratovirt};
                \node[box, below=of dev] (log) {Журнал на хосте};
                \draw[arr] (guest) -- node[right]{запись} (buf);
                \draw[arr] (guest.east) to[bend left=35] node[right]{KICK / FLUSH} (dev.east);
                \draw[arr] (dev) -- node[right]{чтение} (buf);
                \draw[arr] (dev) -- (log);
            \end{tikzpicture}
        \end{column}
    \end{columns}
\end{frame}

\begin{frame}{Виртуальное устройство LogDev}
    \begin{itemize}
        \item Устройство подключается к SysBus Stratovirt и описывается в ACPI как \texttt{STVRT\_LOG}.
        \item Основные MMIO-регистры:
    \end{itemize}

    \vspace{0.2cm}
    \begin{table}
        \small
        \begin{tabular}{llp{0.55\textwidth}}
            \toprule
            Регистр & Доступ & Назначение \\
            \midrule
            \texttt{SYNC\_WR} & запись & синхронный вывод до появления буфера \\
            \texttt{BUFPTR} & чтение/запись & адрес кольцевого буфера в памяти гостя \\
            \texttt{KICK} & запись & уведомление о новых данных \\
            \texttt{FLUSH} & запись & принудительное чтение буфера \\
            \texttt{BUFSIZE} & чтение & размер полезной области буфера \\
            \bottomrule
        \end{tabular}
    \end{table}
\end{frame}

\begin{frame}{MPSC-кольцевой буфер}
    \begin{columns}[T]
        \begin{column}{0.54\textwidth}
            \begin{itemize}
                \item MPSC соответствует модели: много источников в госте и один читатель на хосте.
                \item В \texttt{SHADOW\_W\_POS} хранится мнимая позиция записи и число активных писателей.
                \item \textbf{reserve}: CAS резервирует место и увеличивает счётчик писателей.
                \item \textbf{commit}: CAS уменьшает счётчик; при нуле публикуется \texttt{WritePos}.
                \item Читатель видит только опубликованную область до \texttt{WritePos}.
            \end{itemize}
        \end{column}
        \begin{column}{0.43\textwidth}
            \centering
            \begin{tikzpicture}[
                stage/.style={draw, rounded corners, minimum width=2.8cm, minimum height=0.75cm, align=center},
                arr/.style={-{Latex[length=2mm]}, thick}
            ]
                \node[stage] (r) {reserve\\CAS(\texttt{ShWPos})};
                \node[stage, below=0.45cm of r] (c) {copy\\сообщение};
                \node[stage, below=0.45cm of c] (m) {commit\\CAS + \texttt{WritePos}};
                \draw[arr] (r) -- (c);
                \draw[arr] (c) -- (m);
            \end{tikzpicture}
        \end{column}
    \end{columns}
\end{frame}

\begin{frame}{Поддержка UEFI}
    \begin{itemize}
        \item На стадиях SEC и ранней PEI оперативная память недоступна, поэтому используется \texttt{SYNC\_WR}.
        \item После инициализации RAM PEI-модуль выделяет runtime-память под кольцевой буфер.
        \item Адрес буфера передаётся устройству через \texttt{BUFPTR}.
        \item DXE-драйвер публикует протокол записи и подписывается на \texttt{SetVirtualAddressMap}.
        \item Реализованная \texttt{DebugLib} перенаправляет стандартные \texttt{DEBUG} и \texttt{ASSERT} в LogDev.
    \end{itemize}
\end{frame}

\begin{frame}{Поддержка Windows}
    \begin{itemize}
        \item PnP Manager обнаруживает устройство по ACPI-описанию.
        \item WDM-драйвер получает MMIO-ресурсы, адрес и размер буфера.
        \item Для пользовательского режима предоставлен IOCTL для записи диагностических сообщений.
        \item Для драйверов режима ядра предоставлен интерфейс получения функции записи.
        \item При успешной записи драйвер уведомляет эмулятор через \texttt{KICK}.
    \end{itemize}
\end{frame}

\begin{frame}{Оценка решения}
    \begin{itemize}
        \item Функционально проверяются раннее журналирование UEFI, переход к буферу, несколько источников и обработка переполнения.
        \item Производительность оценивается по задержке доставки и стоимости вызова записи.
        \item Сравниваются три варианта: \texttt{virtio-serial}, LogDev с атомарным резервированием и LogDev с блокировкой.
        \item Важно разделять задержку доставки до хоста и накладные расходы внутри гостевой системы.
    \end{itemize}
\end{frame}

\begin{frame}{Результаты измерений}
    \begin{table}
        \small
        \begin{tabular}{lccc}
            \toprule
            Вариант & Потоки & Средняя доставка, мс & Вызовов/с \\
            \midrule
            \texttt{virtio-serial} & 1--16 & $\approx 70$ & 3.4--4.1 тыс. \\
            LogDev, атомарное резервирование & 1--16 & 336--353 & 8.8--20.1 тыс. \\
            LogDev, блокировка & 1--16 & 799--805 & 12.0--20.4 тыс. \\
            \bottomrule
        \end{tabular}
    \end{table}

    \vspace{0.25cm}
    \begin{itemize}
        \item \texttt{virtio-serial} быстрее доставляет сообщение до хоста.
        \item LogDev имеет меньшую стоимость вызова записи в гостевой системе.
        \item Для виртуализации ключевой показатель --- стоимость вызова записи в гостевой системе; по нему варианты LogDev быстрее \texttt{virtio-serial}.
    \end{itemize}
\end{frame}

\begin{frame}{Результаты работы}
    \begin{itemize}
        \item Разработана архитектура канала диагностики для Stratovirt.
        \item Спроектировано виртуальное устройство LogDev и набор MMIO-регистров.
        \item Обоснован и описан MPSC-кольцевой буфер.
        \item Реализована поддержка UEFI, включая раннюю загрузку и \texttt{DebugLib}.
        \item Реализован Windows-драйвер с интерфейсами для user mode и kernel mode.
        \item Проведено сравнение с \texttt{virtio-serial} и eBPF-подобной организацией буфера.
    \end{itemize}
\end{frame}

\begin{frame}{Заключение}
    \begin{itemize}
        \item Предложенный механизм объединяет раннее журналирование UEFI и быстрый режим через общую память.
        \item MPSC-буфер подходит для нескольких источников диагностических сообщений.
        \item MMIO-регистры используются лишь для синхронного раннего режима, настройки и уведомлений.
        \item Дальнейшая работа: поддержка других ОС и расширение модели безопасности.
    \end{itemize}

    \vspace{0.5cm}
    \centering
    \Large Спасибо за внимание!
\end{frame}

\end{document}
